\section{Conclusion}


Security requires explicit threat models, not vague assertions. This document
provides such a model for the Lux Network, covering four adversary classes, four
protocol layers, three cryptographic assumption families, and their interactions.

The key structural property is defense in depth: any two of the three cryptographic
assumption families (discrete log, lattice, hash) suffice for security. No single
breakthrough---quantum computers, lattice attacks, hash function weaknesses---is
sufficient to compromise the system. Only a simultaneous failure of two independent
mathematical hardness assumptions could break the security guarantees.

The key process property is the machine-checkable invariant: every security claim is
either formally proven, empirically tested, or explicitly marked as an assumption.
This eliminates the most dangerous category of security claims---those that nobody
has checked.

Triple-proof quantum finality achieves the practical synthesis: 248 bytes per epoch,
${\sim}2$--$5$\,ms CPU verification~\cite{lux-benchmarks}, 1\,ms blocks, and security
against classical, quantum, and economic adversaries simultaneously.

%% ========================================================================
%%  BIBLIOGRAPHY
%% ========================================================================
\begin{thebibliography}{20}

\bibitem{boneh2001bls}
D. Boneh, B. Lynn, and H. Shacham, ``Short Signatures from the Weil Pairing,''
\emph{Journal of Cryptology}, vol.~17, no.~4, pp.~297--319, 2004.

\bibitem{fips204}
NIST, ``Module-Lattice-Based Digital Signature Standard (FIPS 204),''
National Institute of Standards and Technology, August 2024.

\bibitem{fips203}
NIST, ``Module-Lattice-Based Key-Encapsulation Mechanism Standard (FIPS 203),''
National Institute of Standards and Technology, August 2024.

\bibitem{fips205}
NIST, ``Stateless Hash-Based Digital Signature Standard (FIPS 205),''
National Institute of Standards and Technology, August 2024.

\bibitem{canetti2021uc}
R. Canetti, R. Gennaro, S. Goldfeder, N. Makriyannis, and U. Peled, ``UC
Non-Interactive, Proactive, Threshold ECDSA with Identifiable Aborts,''
\emph{ACM CCS 2020}, pp.~1769--1787, 2020.

\bibitem{komlo2020frost}
C.~Komlo and I.~Goldberg, ``FROST: Flexible Round-Optimized Schnorr Threshold
Signatures,'' \emph{SAC 2020}, pp.~34--65, Springer, 2020.

\bibitem{rocket2020scalable}
Team Rocket, M. Yin, K. Sekniqi, R. van Renesse, and E.~G. Sirer, ``Scalable and
Probabilistic Leaderless BFT Consensus through Metastability,'' arXiv:1906.08936, 2020.

\bibitem{shor1997}
P.~W. Shor, ``Polynomial-Time Algorithms for Prime Factorization and Discrete
Logarithms on a Quantum Computer,'' \emph{SIAM Journal on Computing}, vol.~26,
no.~5, pp.~1484--1509, 1997.

\bibitem{grover1996}
L.~K. Grover, ``A Fast Quantum Mechanical Algorithm for Database Search,''
\emph{Proceedings of the 28th Annual ACM Symposium on Theory of Computing},
pp.~212--219, 1996.

\bibitem{wood2014ethereum}
G.~Wood, ``Ethereum: A Secure Decentralised Generalised Transaction Ledger,''
Ethereum Project Yellow Paper, 2014.

\bibitem{nakamoto2008bitcoin}
S.~Nakamoto, ``Bitcoin: A Peer-to-Peer Electronic Cash System,'' 2008.

\bibitem{buterin2020combining}
V.~Buterin, D.~Hernandez, T.~Kamphefner, et al., ``Combining GHOST and Casper,''
arXiv:2003.03052, 2020.

\bibitem{lyubashevsky2010ideal}
V.~Lyubashevsky, C.~Peikert, and O.~Regev, ``On Ideal Lattices and Learning with
Errors over Rings,'' \emph{EUROCRYPT 2010}, pp.~1--23, Springer, 2010.

\bibitem{mosca2018}
M.~Mosca, ``Cybersecurity in an Era with Quantum Computers: Will We Be Ready?''
\emph{IEEE Security \& Privacy}, vol.~16, no.~5, pp.~38--41, 2018.

\bibitem{lux-benchmarks}
Lux Industries, ``Quasar Consensus --- Measured Benchmarks,'' 2026. Apple M1 Max, darwin/arm64, Go 1.26.1. \url{https://github.com/luxfi/consensus/blob/main/BENCHMARKS.md}.

\end{thebibliography}

